\documentclass{article}
\usepackage{lmodern}
\usepackage{amsmath}
\usepackage{listings}
\usepackage{fullpage}
\usepackage{parskip}
\usepackage{xcolor}
\lstset{
	basicstyle=\ttfamily,
	columns=fullflexible,
	frame=single,
	breaklines=true,
	postbreak=\mbox{\textcolor{red}{$\hookrightarrow$}\space},
}
\begin{document}
\title{Experiment `p\_6\_2\_c\_replace\_even\_with\_zero\_array' Results}
\author{\today}
\date{}
\maketitle
\textbf{Experiment outcome: }SUCCESS
\\\textbf{Bad responses: }0
\\\textbf{Responses containing}~\texttt{assume}~\textbf{: }0
\\\textbf{Responses containing}~\texttt{expect}~\textbf{: }0
\\\textbf{Resolution attempts: }1
\\\textbf{Hard fails (resolution): }0
\\\textbf{Soft fails (resolution): }0
\\\textbf{Verification attempts: }1
\section*{Problem Specification}
\textbf{Problem name: }p\_6\_2\_c\_replace\_even\_with\_zero\_array
\\\textbf{Natural language statement: }Write a method to replace all even elements in an array with 0.
\\\textbf{Method signature: }p\_6\_2\_c\_replace\_even\_with\_zero\_array(inputs: array<int>) returns (result: array<int>)
\subsection*{Ensures}
\begin{itemize}
\item \texttt{result.Length == inputs.Length}
\item \texttt{forall i :: 0 <= i < inputs.Length ==> (inputs[i] \% 2 == 0 ==> result[i] == 0) \&\& (inputs[i] \% 2 != 0 ==> result[i] == inputs[i])}
\end{itemize}
\clearpage
\section*{GenAI interactions}
Below you will find all interactions between the `user' (program) and the `assistant' (OpenAI).
\subsection*{Program $\rightarrow$ GenAI}
\begin{lstlisting}
You are given the following task to perform in Dafny:

Write a method to replace all even elements in an array with 0.

The signature should be:

p_6_2_c_replace_even_with_zero_array(inputs: array<int>) returns (result: array<int>)

The method should respect the following contract:

ensures result.Length == inputs.Length, ensures forall i :: 0 <= i < inputs.Length ==> (inputs[i] % 2 == 0 ==> result[i] == 0) && (inputs[i] % 2 != 0 ==> result[i] == inputs[i])

Produce and show only the Dafny body of this method, including the curly braces that surround it. Do not show the signature nor contract. You must not use 'assume' nor 'expect' anywhere in your code.
\end{lstlisting}
\subsection*{GenAI $\rightarrow$ Program}
\textbf{System fingerprint: }null
\\\textbf{ID: }chatcmpl-Dn0EibEO82irl0si0f3Y4rAXML6Y1
\\\textbf{Model: }gpt-5.5-2026-04-23
\\\textbf{Created at: }1780571352
\\\textbf{Finish reason: }stop
\\\textbf{Usage: }368 tokens in, and 516 tokens out
\begin{lstlisting}
{
  result := new int[inputs.Length];
  var i := 0;
  while i < inputs.Length
    invariant 0 <= i <= inputs.Length
    invariant result.Length == inputs.Length
    invariant forall j :: 0 <= j < i ==> (inputs[j] % 2 == 0 ==> result[j] == 0) && (inputs[j] % 2 != 0 ==> result[j] == inputs[j])
  {
    if inputs[i] % 2 == 0 {
      result[i] := 0;
    } else {
      result[i] := inputs[i];
    }
    i := i + 1;
  }
}
\end{lstlisting}
\section*{Final Program}
\begin{lstlisting}
method p_6_2_c_replace_even_with_zero_array(inputs: array<int>) returns (result: array<int>)
	ensures result.Length == inputs.Length
	ensures forall i :: 0 <= i < inputs.Length ==> (inputs[i] % 2 == 0 ==> result[i] == 0) && (inputs[i] % 2 != 0 ==> result[i] == inputs[i])
{
  result := new int[inputs.Length];
  var i := 0;
  while i < inputs.Length
    invariant 0 <= i <= inputs.Length
    invariant result.Length == inputs.Length
    invariant forall j :: 0 <= j < i ==> (inputs[j] % 2 == 0 ==> result[j] == 0) && (inputs[j] % 2 != 0 ==> result[j] == inputs[j])
  {
    if inputs[i] % 2 == 0 {
      result[i] := 0;
    } else {
      result[i] := inputs[i];
    }
    i := i + 1;
  }
}
\end{lstlisting}
\section*{Total Token Usage}
\textbf{Input tokens: }368
\\\textbf{Output tokens: }516
\\\textbf{Reasoning tokens: }345
\\\textbf{Sum of `total tokens': }884
\section*{Experiment Timings}
\textbf{Overall Experiment} started at 1780571351985, ended at 1780571382472, lasting 30487ms (30.49 seconds)
\\\textbf{Iteration \#1} started at 1780571351986, ended at 1780571382472, lasting 30486ms (30.49 seconds)
\end{document}
